\documentclass[11pt,aspectratio=169]{beamer}

% ---------- Encoding & Sprache ----------
\usepackage[utf8]{inputenc}
\usepackage[T1]{fontenc}
\usepackage[ngerman]{babel}

\usepackage{graphicx}
\usepackage{listings}
\usepackage{xcolor}

\lstset{
  language=SQL,
  basicstyle=\ttfamily\scriptsize,
  breaklines=true,
  frame=single,
  keywordstyle=\color{blue}\bfseries,
  commentstyle=\color{gray}
}

\usetheme{Madrid}
\setbeamertemplate{navigation symbols}{}

\title[Buchtausch-App]{Datenbankentwurf für eine Buchtausch-App}
\subtitle{Portfolioteil 2 -- Erarbeitungs-/Reflexionsphase}
\author{Tizian Senger}
\institute{DLBDSPBDM01\_D -- Data-Mart-Erstellung in SQL \\ Matrikelnummer: IU14143428}
\date{\today}

\begin{document}

% ---- Folie 1: Titel ----
\begin{frame}
  \titlepage
\end{frame}

% ---- Folie 2: Agenda ----
\begin{frame}{Agenda}
  \tableofcontents
\end{frame}

\section{Überblick}

% ---- Folie 3: Kurzueberblick ueber das Datenbankdesign ----
\begin{frame}{Überblick über das Datenbankdesign}
  % TODO: Kurzer Rueckblick auf das ER-Modell aus Phase 1 (11 Entitaeten, 3 Dreifachbeziehungen).
  \begin{itemize}
    \item Verwendetes DBMS: % TODO
    \item Anzahl Tabellen: % TODO
    \item Normalform: 3. Normalform (3NF)
  \end{itemize}
\end{frame}

\section{Tabellen und SQL-Statements}

% ---- Folien 4-14: Je eine Folie pro Entitaet (Dokumentation, Testfall, Screenshot) ----
% Vorlage fuer jede der 11 Entitaeten aus Phase 1 (main.tex, Abschnitt 4.1):
% Nutzer, Adresse, Buch, Autor, Verlag, Genre, Sprache, Zeitslot, Standort,
% Versandoption, Benachrichtigung

\begin{frame}[fragile]{Tabelle: NUTZER}
  \textbf{CREATE-Statement:}
  \begin{lstlisting}
-- TODO: CREATE TABLE Nutzer (...);
  \end{lstlisting}
  \textbf{Testfall:}
  \begin{lstlisting}
-- TODO: z. B. SELECT * FROM Nutzer WHERE ...;
  \end{lstlisting}
  % TODO: Screenshot des Ergebnisses aus dem DBMS einfuegen
  % \includegraphics[width=\textwidth]{screenshots/nutzer.png}
\end{frame}

\begin{frame}[fragile]{Tabelle: ADRESSE}
  \textbf{CREATE-Statement:}
  \begin{lstlisting}
-- TODO
  \end{lstlisting}
  \textbf{Testfall:}
  \begin{lstlisting}
-- TODO
  \end{lstlisting}
\end{frame}

\begin{frame}[fragile]{Tabelle: BUCH}
  \textbf{CREATE-Statement:}
  \begin{lstlisting}
-- TODO
  \end{lstlisting}
  \textbf{Testfall:}
  \begin{lstlisting}
-- TODO
  \end{lstlisting}
\end{frame}

\begin{frame}[fragile]{Tabelle: AUTOR}
  \textbf{CREATE-Statement:}
  \begin{lstlisting}
-- TODO
  \end{lstlisting}
  \textbf{Testfall:}
  \begin{lstlisting}
-- TODO
  \end{lstlisting}
\end{frame}

\begin{frame}[fragile]{Tabelle: VERLAG}
  \textbf{CREATE-Statement:}
  \begin{lstlisting}
-- TODO
  \end{lstlisting}
  \textbf{Testfall:}
  \begin{lstlisting}
-- TODO
  \end{lstlisting}
\end{frame}

\begin{frame}[fragile]{Tabelle: GENRE}
  \textbf{CREATE-Statement:}
  \begin{lstlisting}
-- TODO
  \end{lstlisting}
  \textbf{Testfall:}
  \begin{lstlisting}
-- TODO
  \end{lstlisting}
\end{frame}

\begin{frame}[fragile]{Tabelle: SPRACHE}
  \textbf{CREATE-Statement:}
  \begin{lstlisting}
-- TODO
  \end{lstlisting}
  \textbf{Testfall:}
  \begin{lstlisting}
-- TODO
  \end{lstlisting}
\end{frame}

\begin{frame}[fragile]{Tabelle: ZEITSLOT}
  \textbf{CREATE-Statement:}
  \begin{lstlisting}
-- TODO
  \end{lstlisting}
  \textbf{Testfall:}
  \begin{lstlisting}
-- TODO
  \end{lstlisting}
\end{frame}

\begin{frame}[fragile]{Tabelle: STANDORT}
  \textbf{CREATE-Statement:}
  \begin{lstlisting}
-- TODO
  \end{lstlisting}
  \textbf{Testfall:}
  \begin{lstlisting}
-- TODO
  \end{lstlisting}
\end{frame}

\begin{frame}[fragile]{Tabelle: VERSANDOPTION}
  \textbf{CREATE-Statement:}
  \begin{lstlisting}
-- TODO
  \end{lstlisting}
  \textbf{Testfall:}
  \begin{lstlisting}
-- TODO
  \end{lstlisting}
\end{frame}

\begin{frame}[fragile]{Tabelle: BENACHRICHTIGUNG}
  \textbf{CREATE-Statement:}
  \begin{lstlisting}
-- TODO
  \end{lstlisting}
  \textbf{Testfall:}
  \begin{lstlisting}
-- TODO
  \end{lstlisting}
\end{frame}

\section{Dreifachbeziehungen}

% ---- Zusaetzliche Folien fuer die 3 Dreifachbeziehungs-Tabellen ----
\begin{frame}[fragile]{Beziehungstabelle: LEIHT\_AUS}
  % Nutzer, Buch, Zeitslot
  \begin{lstlisting}
-- TODO
  \end{lstlisting}
\end{frame}

\begin{frame}[fragile]{Beziehungstabelle: BEWERTET}
  % Nutzer (Bewertende:r), Buch, Nutzer (Eigentuemer:in)
  \begin{lstlisting}
-- TODO
  \end{lstlisting}
\end{frame}

\begin{frame}[fragile]{Beziehungstabelle: VERFÜGBAR\_AN}
  % Buch, Zeitslot, Standort
  \begin{lstlisting}
-- TODO
  \end{lstlisting}
\end{frame}

\section{Zusammenfassung}

% ---- Vorletzte Folie: Kurze Zusammenfassung der Implementierung (mind. 0,5 Seiten Inhalt) ----
\begin{frame}{Zusammenfassung der Implementierung}
  % TODO: mind. eine halbe Seite Text/Stichpunkte gemaess Aufgabenstellung.
  % Siehe auch main.tex (KZ) in diesem Ordner fuer die ausformulierte Kurzzusammenfassung.
\end{frame}

% ---- Letzte Folie: Fazit / Ausblick ----
\begin{frame}{Fazit und Ausblick}
  % TODO: Herausforderungen, offene Punkte fuer die Finalisierungsphase.
\end{frame}

\end{document}
